\documentclass[10pt,a4paper]{amsart}
\usepackage[margin=19mm]{geometry}
\usepackage{amsmath,amssymb,mathtools}
\usepackage[scr=rsfs,cal=euler]{mathalfa}
\usepackage{xfrac}
\usepackage[unicode,hidelinks,bookmarks=false]{hyperref}
\newcommand{\ie}{i.e.\ }
\newcommand{\cf}{cf.\ }
\newcommand{\resp}{respectively\ }
\newcommand{\Subsection}{\subsection}
\providecommand{\nobreakdash}{\nobreak\hbox{-}}
\setlength{\emergencystretch}{2em}
\multlinegap0pt
\usepackage{amssymb}
\usepackage{enumerate}
\usepackage{xcolor}
\usepackage{float}
\usepackage[shortlabels]{enumitem}
\usepackage{tikz}
\newtheorem{theorem}{Theorem}
\title{The sad life of lattice triangles}
\author{Christian Aebi, Grant Cairns}
\date{Source: arXiv:2602.23392v1 (CC BY 4.0)}
\begin{document}
\maketitle

\noindent\emph{Source and licence.} The sad life of lattice triangles. Version arXiv:2602.23392v1 (CC BY 4.0), distributed by arXiv under the Creative Commons Attribution 4.0 International licence, http://creativecommons.org/licenses/by/4.0/, as recorded in the arXiv licence metadata for this exact version and shown on the abstract page https://arxiv.org/abs/2602.23392v1. Full licence notice: \texttt{licenses/arxiv-2602.23392-CC-BY-4.0.txt}.\par\smallskip

\noindent\textit{Editorial scope.} Counted unit: the article's theorem theorem (source0.tex lines 218--220). The article's complete original proof of it is delivered with it (source0.tex lines 222--234, one \texttt{proof} environment of 13 lines). No mathematics is written, repaired or re-derived: the statement and the proof are the article's own lines, in the article's own notation, and no retained line is substituted or re-flowed.  Retained uncounted as the source-local context the proof needs: lines 169--171.  Nothing was omitted from any retained span, so the package carries no omission record.  No retained line contains a citation, so the wrapper carries no bibliography and no bibliography entry.  No retained line contains a figure environment, an image-inclusion command or drawing code, so no image is delivered and the package declares no asset dependency.  Editorial formatting only, in addition: an AMS \texttt{amsart} preamble that restates the article's own declarations and macros used by the retained lines, namely the environments \texttt{theorem} on separate counters, together with the standard packages those lines need.  The retained environments print in this document as Theorem 1 in order of appearance, whereas the article prints the same items with the numbering of its own full document.  The complete native source of the article as distributed in the arXiv e-print for this exact version is bundled unchanged as \texttt{sources/195396/source0.tex}.
\begin{equation}\label{E:Omeganum}
v_1v_2(\overline{v_2}-\overline{v_1})=(x_2^2+y_2^2)(x_1,y_1) - (x_1^2+y_1^2)(x_2,y_2).
\end{equation}

\begin{theorem}
If for the lattice triangle with vertices $O,(x_1,y_1),(x_2,y_2)$, the circumcenter  and the centroid are both lattice points, then $\gcd(x_1,y_1,x_2,y_2)$ is a multiple of $3$.
\end{theorem}

\begin{proof}
Naturally enough, the idea is  to work simply modulo $3$. If the centroid $G$ is lattice point, then 
\begin{equation}\label{E:xxyy}
x_1\equiv -x_2\pmod3\qquad\text{and}\qquad y_1\equiv -y_2\pmod3.
\end{equation}
Now looking at the formula for the circumcenter $F$,  and using the above relations,  we see the denominator of $F$ is $2(x_1 y_2-y_1 x_2) \equiv 0\pmod3$. As $F$ is a lattice point, the numerator of $F$ must then also be divisible by $3$.
From \eqref{E:Omeganum} and \eqref{E:xxyy}, the numerator, when evaluated modulo $3$,
 gives
\[
(x_2^2+y_2^2)(x_1,y_1) - (x_1^2+y_1^2)(x_2,y_2)\equiv 2(x_1^2+y_1^2)(x_1,y_1)\pmod3.
\]
Since squares are congruent to either $0$ or $1$ modulo $3$ according to whether or not the number is a multiple of $3$, it follows that if the sum of two squares is divisible by $3$ then both terms are multiples of $3$. So the above congruence implies necessarily that $x_1$ and $y_1$ are both divisible by $3$, and by  \eqref{E:xxyy} the same is true for $x_2$ and $y_2$, i.e., $\gcd(x_1,y_1,x_2,y_2)$ is a multiple of $3$.
\end{proof}

\end{document}
